% !TEX TS-program = xelatex
% !TEX encoding = UTF-8 Unicode
% !Mode:: "TeX:UTF-8"

\documentclass{resume}
\usepackage{zh_CN-Adobefonts_external} % Simplified Chinese Support using external fonts (./fonts/zh_CN-Adobe/)
% \usepackage{NotoSansSC_external}
% \usepackage{NotoSerifCJKsc_external}
% \usepackage{zh_CN-Adobefonts_internal} % Simplified Chinese Support using system fonts
\usepackage{linespacing_fix} % disable extra space before next section
\usepackage{cite}

\begin{document}
\pagenumbering{gobble} % suppress displaying page number

\name{程勇 - C++ 开发工程师 - 随时入职}

\basicInfo{
  \email{xhcoding@foxmail.com} \textperiodcentered\
  \phone{(+86) 18360868781} \textperiodcentered\
  \faBirthdayCake{1998-10} \textperiodcentered\
  \faGithubSquare{https://github.com/xhcoding}}

\section{\faCogs\ 个人能力}
\begin{itemize}
  \item 编程语言：熟悉 C/C++ ，Python ，较为熟悉 JavaScript ，Bash ，Dart ，Rust ，ELisp
  \item 设计能力：熟悉常用设计模式，面向对象设计，能够进行需求分析及功能设计
  \item 开发能力：熟悉常用的数据结构及算法，多线程及多进程编程，高性能网络编程，尤其熟悉 TCP/IP 协议栈
  \item 常用工具：熟悉 Windows/Linux 下常用的开发调试工具，如 Git ，WireShark ，WinDbg ，GDB ，Perf ，Emacs 等
  \item 个人优势：大厂经验，熟悉成熟的研发流程；工程能力强，学习能力强，能够快速适应大型复杂项目的开发；有技术热情，攻关能力强
\end{itemize}

\section{\faGraduationCap\ 教育经历}
\datedsubsection{\textbf{河海大学}, 南京}{2015 年 9 月 -- 2019 年 6 月}
\textit{学士}\ 计算机科学与技术

\section{\faWrench\ 工作经历}
\datedsubsection{\textbf{华为技术有限公司(OD)}, 东莞}{2023 年 5 月 -- 2025 年 10 月}
从事网络(TCP/IP)协议栈相关工作，在团队中担任 MDE ：
\begin{itemize}
  \item 个人职责: 项目原型技术验证及性能优化，达成关键指标，消除项目商用风险，项目功能设计及核心模块开发等
  \item 团队职责: 构建工具，编辑器插件及测试底座的开发，提升团队开发及测试效率，日常代码检视，设计评审等
\end{itemize}


\datedsubsection{\textbf{视源电子科技股份有限公司}, 广州}{2019 年 7 月 -- 2022 年 3 月}
基于 QT/QML 进行 PC 客户端开发；使用 C++ 将系统功能封装为 JavaScript API ；设备网关中间件开发

\section{\faHeartO\ 荣誉}
\datedline{\textbf{PDU 年度优秀个人}（绩效前 1\%）}{2024 年 12 月}

\section{\faGroup\ 项目经历}
\datedsubsection{\textbf{ 网络协议栈 }}{2023年 11 月 -- 至今}
\textbf{ 背景： }
\begin{itemize}
  \item C 语言实现 ETH ，ARP ，IP/IP6 ，TCP/UDP ，DHCP ，DNS 等协议，用于嵌入式小型化场景，用户态高性能场景及鸿蒙内核场景
\end{itemize}
\textbf{ 职责： }
\begin{itemize}
  \item 性能优化：根据客户场景设计性能模型，实现性能工程，基线性能数据，通过各种工具分析性能瓶颈点后针对性优化，最终达成客户性能目标。
    在嵌入式小型化场景，协议栈内存占用及 iPerf 打流性能均优于业界的 LwIP ，提升 xx\% ；用户态高性能场景 Redis 业务及 iPerf 打流性能
    对比 Linux 内核提升 xxx\% ；鸿蒙内核场景 iPerf 打流性能对比 Linux 内核提升 xx\%
  \item 设计开发：负责协议栈部分需求分析及功能设计，编写功能设计文档；
    负责核心模块的开发和单元测试，包括 TCP ，IP 等协议实现，Socket 及 Epoll 事件接口等；
    负责构建工程及一些测试工程底座的开发。
    开发效率高，质量好，多次荣获团队迭代之星
  \item 问题攻关：负责疑难问题的定位攻关，多次高效的处理紧急任务
  \item 团队赋能：多次进行业务培训，开发编辑器插件，优化用到的上游库及框架，提升开发测试效率
\end{itemize}

\textbf{ 相关技术： } C/C++ ，Python ，DPDK ，Linux 内核网络栈，TCP/IP 相关的网络协议

\datedsubsection{\textbf{ HTTP1/2 客户端/服务器/代理 }}{2023年 5 月 -- 2023 年 11 月}
\textbf{ 背景： }
\begin{itemize}
  \item 预研项目，高性能 HTTP1/2 的客户端，服务器及代理，用于替代 nginx ，目标是性能优于 nginx
\end{itemize}
\textbf{ 职责： }
\begin{itemize}
  \item 技术预研：在已有 HTTP1 的基础上，集成公司其它部门用 Rust 写的 HTTP2 组件，实现 HTTP1/2 客户端，服务端及代理功能，并进行性能优化，
    最终在年底达成预研目标，性能超过 nginx 。
\end{itemize}

\textbf{ 相关技术： } C/C++ ，HTTP1 和 HTTP2 协议

\datedsubsection{\textbf{ 设备接口网关 }}{2021年 1 月 -- 2022 年 3 月}
\textbf{ 背景： }
\begin{itemize}
  \item 以接口形式向应用提供设备的硬件能力，包括摄像头，阵列麦克风，MCU ，蓝牙 BLE 外设等。
\end{itemize}
\textbf{ 职责： }
\begin{itemize}
  \item 实现多种通信链路及协议间的通信，包括 HTTP ，MQTT ，HID ，蓝牙 BLE 等
  \item 实现各种语言接入网关的 SDK ，方便各类应用使用
  \item 使用 flutter 开发可视化的测试工具，提高测试效率
\end{itemize}

\textbf{ 相关技术： } C/C++ ，Boost ，Asio ，HTTP/MQTT/MCU/HID 通信等

\datedsubsection{\textbf{ node-localsdk }}{2019年 7 月 -- 2022 年 3 月}
\textbf{ 背景： }
\begin{itemize}
  \item 使用 C++ 将 Windows 系统和 Linux 系统的一些功能封装成 JavaScript 接口，供 Electron 应用调用
\end{itemize}
\textbf{ 职责： }
\begin{itemize}
  \item 根据需要设计 API 并用 C++ 实现，导出为 JavaScript 接口
  \item 项目完全满足前端工程化的要求，自动生成文档并发布，完善的单元测试，保证提供的 SDK 质量
\end{itemize}

\textbf{ 相关技术： } C++, NodeJS, Catch2, Mocha, N-API, Win32 API, Linux API 等


\datedsubsection{\textbf{ 系统升级应用 }}{2019年 12 月 -- 2022 年 3 月}
\textbf{ 背景： }
\begin{itemize}
  \item 基于 QT/QML 开发的 PC 端应用程序，提供各种设备（摄像头，麦克风）固件的自动或手动升级功能
\end{itemize}
\textbf{ 职责： }
\begin{itemize}
  \item 参加产品，交互，视觉评审，根据最终需求及视觉交互稿完成开发和上线
  \item 持续优化，让项目具有良好的可扩展性，解决各种 BUG, 深入 QT 源码解决遇到的问题
\end{itemize}

\textbf{ 相关技术： } C++, QT5, Qt Quick, Websocket 通信等


\end{document}
